\chapter{Học Tủ Lệch Đề}
\chapterintrobi{Đoán trúng thì vui, đoán lệch thì rất thơ}{Study the wrong topic}{Đây là kiểu mạo hiểm được nhiều học sinh yêu thích dù xác suất phản đòn khá cao.}

\begin{lucbatbox}{Lục bát: Tủ không mở đúng ngăn}
Em ôn đúng một vùng thôi\
Tin rằng đề sẽ rơi rơi chỗ này\
Ai ngờ đời chẳng chiều tay\
Mở ra câu khác đứng đầy trang kia\
Tủ em khóa rất say mê\
Nhưng chìa sai ổ nên về tay không
\end{lucbatbox}

\begin{tutuyetbox}{Thất ngôn tứ tuyệt: Niềm tin lệch hướng}
Em học rất kỹ mấy phần tin\
Đề ra chỗ khác, dạ lặng thinh\
Nếu ôn rộng chút thay vì đặt cược\
Bài đâu cần hóa cuộc hành trình
\end{tutuyetbox}

\begin{ghichu}
Học tủ là cách mua vé số bằng thời gian ôn tập. Trúng thì vui, lệch thì đau rất rõ.
\end{ghichu}

\begin{englishnote}
\textbf{study the wrong topic}: học lệch đề.\par
\textbf{guess the exam}: đoán đề.\par
\textbf{I studied the wrong part.}: Em học lệch phần rồi.
\end{englishnote}